\documentclass[a4paper,12pt]{scrartcl}

% Layout
\usepackage[a4paper, left=2.5cm, right=2cm, top=2cm, bottom=3.2cm]{geometry} % kleinere Ränder
\usepackage{parskip}

% Umlaute in der Datei erlauben, auf Deutsch umstellen
\usepackage{fontspec}
\usepackage{lmodern}
\usepackage[english, ngerman]{babel}

% Mathesymbole und Ähnliches
\usepackage{amsmath}
\usepackage{amsthm}
\usepackage{mathtools}
\usepackage{amssymb}
\usepackage{microtype}
\usepackage{stmaryrd}
\usepackage{MnSymbol}
\usepackage{bussproofs}		% Sequenzenkalkül
\usepackage{xcolor}

% Grafiken und PDFs einfügen
\usepackage{graphicx}
\usepackage{pdfpages}
\usepackage{float}

% PDF-Tools
\usepackage[hidelinks, unicode]{hyperref}

% Abbildungen
\usepackage{tikz}
\usetikzlibrary{arrows,calc,positioning}
\tikzset{ 
   state/.style={circle,fill,inner sep=2pt,label=below:#1\strut},
   arrow/.style={->, shorten <=3pt, shorten >=3pt, >=stealth',thick},
   knoten/.style={draw, circle, minimum size=8mm}} 
   
% Infoboxen 
\definecolor{orange}{RGB}{230,159,0}
\definecolor{darkgray}{RGB}{70,90,70}

\usepackage{tcolorbox}


% Reelle, Natürliche, Ganze, Rationale Zahlen
\newcommand{\R}{\ensuremath{\mathbb{R}}}
\newcommand{\N}{\ensuremath{\mathbb{N}}}
\newcommand{\Z}{\ensuremath{\mathbb{Z}}}
\newcommand{\Q}{\ensuremath{\mathbb{Q}}}

% Fraktur für Strukturen
\newcommand{\A}{\ensuremath{\mathfrak A}}
\newcommand{\B}{\ensuremath{\mathfrak B}}
\newcommand{\I}{\ensuremath{\mathfrak I}}

% Körper
\newcommand{\K}{\ensuremath{\mathcal{K}}}
%\newcommand{\G}{\ensuremath{\mathcal{G}}}

% Makros für logische Operatoren
\newcommand{\xor}{\ensuremath{\oplus}} % exklusives oder
\newcommand{\impl}{\ensuremath{\rightarrow}} % logische Implikation
\newcommand{\Lor}{\ensuremath{\bigvee}}
\newcommand{\Land}{\ensuremath{\bigwedge}}

% Meistens ist \varphi schöner als \phi, genauso bei \theta
\renewcommand{\phi}{\varphi}
\renewcommand{\theta}{\vartheta}
\renewcommand{\epsilon}{\varepsilon}

\newcommand*{\todo}{\textbf{\textcolor{red}{TODO}}}		% big, red, apparent TODO

% SEQUENZENKALKÜL
\usepackage{bussproofs}
\newcommand{\yields}{\ensuremath{\mathrel{\Rightarrow}}}

% Aufzählungen anpassen (alternativ: \arabic, \alph)
\renewcommand{\labelenumi}{(\roman{enumi})}

% BITTE NICHT ÄNDERN: interne Kommandos für die Informationen (Blattnummer, Gruppe, ...)
\newcounter{sheet}
\newcommand{\printsheet}{\arabic{sheet}}

\newcommand{\printgroup}{?}
\newcommand{\group}[1]{%
\renewcommand{\printgroup}{#1}%
}

\newcommand{\printmembers}{}
\newcommand{\printmember}[3]{{#2} {#3} & {#1} \\}

% Für Vektoren
\newcommand{\myvec}[1]{\ensuremath{\begin{pmatrix}#1\end{pmatrix}}}
\def\spvec#1{\left(\vcenter{\halign{\hfil$##$\hfil\cr \spvecA#1;;}}\right)}
\def\spvecA#1;{\if;#1;\else #1\cr \expandafter \spvecA \fi}

\newcommand{\mal}{\cdot}

\newtheorem{proposition}{Proposition}[sheet]

\newcommand{\pdfmembers}{}
\newcommand{\pdfmember}[3]{{#1} {#3}, {#2}; }

\newcommand{\member}[3]{%
\expandafter\renewcommand\expandafter\printmembers\expandafter{\printmembers\printmember{#1}{#2}{#3}}%
\expandafter\renewcommand\expandafter\pdfmembers\expandafter{\pdfmembers\pdfmember{#1}{#2}{#3}}%
}

\AtBeginDocument{\hypersetup{
    pdftitle = {Übungsblatt \printsheet},
    pdfauthor = {Abgabegruppe \printgroup: \pdfmembers},
    pdfsubject = {Online Algorithms}
}}

% <-------- HIER alle Informationen eintragen ========================
\setcounter{sheet}{4} % Nummer des Blatts

% alle Gruppenmitglieder in der Form \member{Matrikelnummer}{Vorname}{Nachname}
\member{5481114}{Danje}{Petersen}

\begin{document}

% Kopfzeile
{\raggedright
\begin{tabular}{l}
    Online Algorithms \\
    SS 2024 \\
\end{tabular}}
\hfill
{\Large Sheet 0\printsheet}
\hfill
\begin{tabular}{l l}
\today{} \\
\printmembers
\end{tabular}
\hrule


% <-------- HIER beginnt die Lösung ========================

\section*{Exercise \printsheet.1}
\subsection*{(a)}
In the worst case, we almost reach the pasture and travel $D-\epsilon$. After that the following steps need to occur:
\begin{itemize}
\item The same distance needs to be travelled back to the start 
\item The distance $2 \mal (D-\epsilon)$ is travelled in the wrong direction and back
\item We travel the distance $D$ into the correct direction and find the pasture.
\end{itemize} 
This amounts to a total distance of $D_\epsilon + D -\epsilon + 2 \mal (D-\epsilon) + D = 7D - 6\epsilon$ being travelled shortly before and after we almost reached the pasture. Since we can choose $\epsilon$ arbitrarily small while being $>0$ we can ignore it for further argumentation. \\

For the time before we travelled $D-\epsilon$ for the first time, we must have walked the distances $\frac{D-\epsilon}{2^i}$ for $i \in \{1,\dots \log_2(D)\}$ twice until we get back to the first step of moving one unit.
The remaining distances can summed up in this term:
\[\sum_{i=1}^{\log_2(D)} \frac{2D}{2^{i}} = \sum_{i=1}^{\log_2(D)} \frac{D}{2^{i-1}} < 2D\]

This makes the total distance being travelled $7D + \sum_{i=1}^{\log_2(D)} \frac{D}{2^{i-1}} < 9D$. Since OPT would simply travel to the pasture $c_{OPT}=D$ and therefore \[c = \frac{c_{ALG}}{c_{OPT}} = \frac{9D}{D} = 9\]

Since $D$ could be arbitrarily high, we cannot set a lower competitive ratio than $9$.
The higher $D$ is, the closer to $9$ would the factor between the strategy and OPT become.  

\subsection*{(b)}
The restriction $D \geq 1$ is necessary because else the adversary could make $D$ as small as possible and place it in the direction where we are not travelling in the first move. Since the first move is always one unit long and $D$ can be arbitrarily small the competitive factor would be unbounded.
\[c = \frac{c_{ALG}}{c_{OPT}} = \frac{1}{D} = \infty \qquad \mid \text{for } \lim_{D \rightarrow 0}\]



\section*{Exercise \printsheet.2}



\section*{Exercise \printsheet.3}



\section*{Exercise \printsheet.4}





\end{document}
